\begin{alist}
\item Το τριώνυμο $\omega^2 + 4\omega - 12$ έχει ρίζες τις $\omega = 2$ και $\omega = -6$. Για $\omega \neq 2$ και $\omega \neq -6$ έχουμε ισοδύναμα
\[\frac{\omega^3 - 8}{\omega^2 + 4\omega - 12} > 0 \Leftrightarrow \frac{(\omega - 2)(\omega^2 + 2\omega + 4)}{(\omega - 2)(\omega + 6)} > 0 \Leftrightarrow \frac{\omega^2 + 2\omega + 4}{\omega + 6} > 0\]
και επειδή το τριώνυμο $\omega^2 + 2\omega + 4$ έχει αρνητική διακρίνουσα θα είναι για κάθε $\omega \in \mathbb{R}$ ομόσημο του συντελεστή του $\omega^2$, δηλαδή θετικό, έχουμε τελικά ότι
\[\omega + 6 > 0 \Leftrightarrow \omega > -6.\]
Το σύνολο των λύσεων της ανίσωσης $\dfrac{\omega^3 - 8}{\omega^2 + 4\omega - 12} > 0$ είναι το $(-6, 2) \cup (2, +\infty)$.

\item Η συνάρτηση $f$ ορίζεται για τις τιμές του $x \in \mathbb{R}$ για τις οποίες ισχύει
\[\frac{e^{3x} - 8}{e^{2x} + 4e^x - 12} > 0.\]
Αν θέσουμε $e^x = \omega$ η τελευταία ανίσωση γίνεται
\[\frac{\omega^3 - 8}{\omega^2 + 4\omega - 12} > 0\]
που όπως δείξαμε στο $\alpha$) αληθεύει για κάθε $\omega \in (-6, 2) \cup (2, +\infty)$. \\
Συνεπώς θα πρέπει $e^x \neq 2 \Leftrightarrow x \neq \ln 2$ και $e^x > -6$ που ισχύει. Τελικά το πεδίο ορισμού της $f$ είναι το $\mathbb{R} - \{\ln 2\}$.

\item Οι τετμημένες των σημείων τομής της γραφικής παράστασης της $f$ με τον άξονα $xx'$ είναι οι λύσεις της εξίσωσης $f(x) = 0$ με $x \neq \ln 2$. Είναι:
\begin{align*}
f(x) = 0 &\Leftrightarrow \ln\frac{e^{3x} - 8}{e^{2x} + 4e^x - 12} = \ln 1 \Leftrightarrow \\
\frac{e^{3x} - 8}{e^{2x} + 4e^x - 12} &= 1 \Leftrightarrow \\
e^{3x} - 8 &= e^{2x} + 4e^x - 12 \Leftrightarrow \\
e^{3x} - e^{2x} - 4e^x + 4 &= 0 \Leftrightarrow \\
e^{2x}(e^x - 1) - 4(e^x - 1) &= 0 \Leftrightarrow \\
(e^x - 1)(e^{2x} - 4) &= 0 \Leftrightarrow \\
(e^x - 1)(e^x - 2)(e^x + 2) &= 0.
\end{align*}
Συνεπώς θα πρέπει $e^x - 1 = 0 \Leftrightarrow e^x = 1 \Leftrightarrow x = 0$ που είναι δεκτή ή $e^x - 2 = 0 \Leftrightarrow e^x = 2 \Leftrightarrow x = \ln 2$ που απορρίπτεται ή $e^x + 2 = 0 \Leftrightarrow e^x = -2$ που είναι αδύνατη.
\end{alist}
